% Mathématiques 4e — cours V3
% Translation, agrandissement et réduction

\section{Objectifs}

\begin{itemize}
\item Reconnaître et construire l'image d'une figure par translation.
\item Décrire une translation à l'aide de deux points ou d'un vecteur de déplacement.
\item Utiliser les propriétés de conservation d'une translation.
\item Comprendre un agrandissement et une réduction comme transformations homothétiques simples.
\item Calculer une longueur après agrandissement ou réduction.
\item Déterminer un coefficient d'agrandissement ou de réduction.
\item Étudier l'effet d'un coefficient sur les longueurs, les aires et les volumes.
\item Résoudre des problèmes de plans, maquettes, photographies et solides semblables.
\end{itemize}

\section{1. Déplacer sans déformer}

Une translation déplace tous les points d'une figure de la même manière.

Chaque point se déplace :
\begin{itemize}
\item dans la même direction ;
\item dans le même sens ;
\item de la même longueur.
\end{itemize}

\coursefigure{/images/mathematiques/4eme/translation.svg}{Deux triangles identiques reliés par des déplacements parallèles de même longueur et de même sens.}{Une translation déplace une figure sans modifier sa forme ni ses dimensions.}

\section{2. Image d'un point}

Si le point $A$ a pour image $A'$ par une translation, le déplacement de $A$ vers $A'$ caractérise la transformation.

Si $B$ est un autre point de la figure et $B'$ son image, alors les déplacements :
\[
\overrightarrow{AA'}
\]
et :
\[
\overrightarrow{BB'}
\]
sont identiques.

On peut écrire :
\[
\boxed{\overrightarrow{AA'}=\overrightarrow{BB'}}.
\]

\section{3. Construire l'image d'un point}

Pour construire $B'$, image de $B$ par la translation qui transforme $A$ en $A'$ :
\begin{enumerate}
\item tracer par $B$ une droite parallèle à $(AA')$ ;
\item reporter la longueur $AA'$ ;
\item respecter le sens du déplacement ;
\item placer $B'$.
\end{enumerate}

Le quadrilatère $ABB'A'$ est alors un parallélogramme lorsque les points ne sont pas alignés.

\section{4. Propriétés conservées}

Une translation conserve :
\begin{itemize}
\item les longueurs ;
\item les angles ;
\item l'alignement ;
\item le parallélisme ;
\item les aires.
\end{itemize}

Ainsi, si :
\[
AB=5\ \text{cm},
\]
alors :
\[
A'B'=5\ \text{cm}.
\]

\section{5. Image d'un segment}

L'image du segment $[AB]$ est le segment $[A'B']$.

On a :
\[
AB=A'B'.
\]

De plus :
\[
(AB)\parallel(A'B').
\]

La translation conserve donc à la fois la longueur et la direction du segment.

\section{6. Image d'un cercle}

Un cercle de centre $O$ et de rayon $r$ devient un cercle de centre $O'$ et de même rayon :
\[
r.
\]

La translation déplace le centre, mais ne modifie pas le rayon.

\section{7. Translation dans un repère}

Dans un repère, une translation peut être décrite par un déplacement horizontal et vertical.

Par exemple :
\[
\Delta x=4,
\qquad
\Delta y=-2.
\]

Le point :
\[
A(3;5)
\]
a alors pour image :
\[
A'(7;3).
\]

En effet :
\[
3+4=7,
\]
et :
\[
5-2=3.
\]

\section{8. Reconnaître une translation}

Pour reconnaître une translation entre deux figures, il faut vérifier que les points correspondants ont subi le même déplacement.

Une rotation ou une symétrie peut conserver les longueurs elle aussi, mais le déplacement des points n'est pas identique.

\begin{attention}
Deux figures de même forme et de même taille ne sont pas forcément images l'une de l'autre par translation.
\end{attention}

\section{9. Agrandissement}

Un agrandissement multiplie toutes les longueurs par un même nombre :
\[
k>1.
\]

Ce nombre est le coefficient d'agrandissement.

Si :
\[
AB=6\ \text{cm}
\]
et :
\[
k=1{,}5,
\]
alors :
\[
A'B'=1{,}5\times6=9\ \text{cm}.
\]

\section{10. Réduction}

Une réduction multiplie toutes les longueurs par un coefficient :
\[
0<k<1.
\]

Par exemple :
\[
k=0{,}4.
\]

Une longueur de :
\[
15\ \text{cm}
\]
devient :
\[
15\times0{,}4=6\ \text{cm}.
\]

\section{11. Déterminer le coefficient}

Si une longueur :
\[
8\ \text{cm}
\]
devient :
\[
14\ \text{cm},
\]
le coefficient est :
\[
k=\dfrac{14}{8}=1{,}75.
\]

Il s'agit d'un agrandissement car :
\[
k>1.
\]

\section{12. Conserver les angles}

Un agrandissement ou une réduction conserve les angles.

Si un triangle possède un angle de :
\[
52^\circ,
\]
l'angle correspondant dans la figure agrandie mesure encore :
\[
52^\circ.
\]

La figure conserve donc sa forme.

\section{13. Longueurs correspondantes}

Si deux figures sont obtenues par agrandissement ou réduction de coefficient $k$, alors :
\[
\boxed{
\text{longueur image}=k\times\text{longueur initiale}
}.
\]

Tous les rapports de longueurs correspondantes valent :
\[
k.
\]

\section{14. Effet sur les aires}

Une aire est obtenue à partir de deux dimensions.

Si chaque longueur est multipliée par :
\[
k,
\]
alors l'aire est multipliée par :
\[
\boxed{k^2}.
\]

\coursefigure{/images/mathematiques/4eme/agrandissement-effets.svg}{Résumé des effets d'un coefficient $k$ sur les longueurs, les aires et les volumes.}{Un même coefficient $k$ agit différemment selon la dimension de la grandeur : $k$ sur les longueurs, $k^2$ sur les aires et $k^3$ sur les volumes.}

\section{15. Exemple sur une aire}

Un carré de côté :
\[
4\ \text{cm}
\]
a pour aire :
\[
16\ \text{cm}^2.
\]

On l'agrandit avec :
\[
k=3.
\]

Le nouveau côté vaut :
\[
12\ \text{cm}.
\]

La nouvelle aire vaut :
\[
144\ \text{cm}^2.
\]

On vérifie :
\[
144=16\times3^2.
\]

\section{16. Effet sur les volumes}

Un volume dépend de trois dimensions.

Si toutes les longueurs sont multipliées par :
\[
k,
\]
alors le volume est multiplié par :
\[
\boxed{k^3}.
\]

\section{17. Exemple sur un volume}

Un cube d'arête :
\[
2\ \text{cm}
\]
a pour volume :
\[
2^3=8\ \text{cm}^3.
\]

On agrandit avec :
\[
k=2{,}5.
\]

La nouvelle arête vaut :
\[
5\ \text{cm}.
\]

Le nouveau volume vaut :
\[
5^3=125\ \text{cm}^3.
\]

On vérifie :
\[
125=8\times2{,}5^3.
\]

\section{18. Retrouver un coefficient à partir d'une aire}

Une figure d'aire :
\[
20\ \text{cm}^2
\]
devient une figure semblable d'aire :
\[
45\ \text{cm}^2.
\]

Le facteur sur les aires vaut :
\[
\dfrac{45}{20}=2{,}25.
\]

Donc :
\[
k^2=2{,}25.
\]

Comme $k>0$ :
\[
k=1{,}5.
\]

\section{19. Maquette}

Une maquette est réalisée à l'échelle :
\[
1:50.
\]

Le coefficient entre la réalité et la maquette vaut :
\[
k=\dfrac1{50}.
\]

Une hauteur réelle de :
\[
8\ \text{m}=800\ \text{cm}
\]
devient :
\[
800\times\dfrac1{50}=16\ \text{cm}.
\]

\section{20. Agrandissement photographique}

Une photographie mesure :
\[
10\ \text{cm}\times15\ \text{cm}.
\]

On veut une largeur de :
\[
24\ \text{cm}.
\]

Le coefficient vaut :
\[
k=\dfrac{24}{15}=1{,}6.
\]

La hauteur devient :
\[
10\times1{,}6=16\ \text{cm}.
\]

Le nouveau format est :
\[
\boxed{16\ \text{cm}\times24\ \text{cm}}.
\]

\section{21. Aire d'une photographie}

L'aire initiale vaut :
\[
10\times15=150\ \text{cm}^2.
\]

Avec :
\[
k=1{,}6,
\]
l'aire est multipliée par :
\[
1{,}6^2=2{,}56.
\]

Donc :
\[
150\times2{,}56=384\ \text{cm}^2.
\]

On retrouve :
\[
16\times24=384\ \text{cm}^2.
\]

\section{22. Ne pas confondre translation et agrandissement}

Une translation conserve :
\[
\text{les longueurs}.
\]

Un agrandissement de coefficient :
\[
k\neq1
\]
modifie les longueurs.

Ainsi :
\[
AB=A'B'
\]
pour une translation, tandis que :
\[
A'B'=kAB
\]
pour un agrandissement ou une réduction.

\section{23. Méthode : translation}

\begin{methode}
\begin{enumerate}
\item identifier le déplacement donné ;
\item repérer direction, sens et longueur ;
\item appliquer exactement ce déplacement à chaque point utile ;
\item reconstruire la figure image ;
\item vérifier les longueurs et le parallélisme.
\end{enumerate}
\end{methode}

\section{24. Méthode : agrandissement ou réduction}

\begin{methode}
\begin{enumerate}
\item identifier deux longueurs correspondantes ;
\item calculer le coefficient $k$ ;
\item décider s'il s'agit d'un agrandissement ou d'une réduction ;
\item multiplier toutes les longueurs correspondantes par $k$ ;
\item utiliser $k^2$ pour les aires et $k^3$ pour les volumes ;
\item contrôler la cohérence du résultat.
\end{enumerate}
\end{methode}

\section{25. Erreurs fréquentes}

\begin{itemize}
\item Déplacer différents points avec des directions différentes dans une translation.
\item Modifier les longueurs lors d'une translation.
\item Appliquer $k$ aux aires au lieu de $k^2$.
\item Appliquer $k^2$ aux volumes au lieu de $k^3$.
\item Confondre une réduction de coefficient $0{,}5$ avec une diminution de $0{,}5$ unité.
\item Utiliser des longueurs non correspondantes pour déterminer $k$.
\item Oublier de convertir les unités dans un problème d'échelle.
\end{itemize}

\section{À retenir}

\begin{aretenir}
Une translation déplace chaque point selon le même déplacement et conserve :
\[
\boxed{\text{longueurs, angles, alignement, parallélisme et aires}}.
\]

Dans un agrandissement ou une réduction de coefficient $k$ :
\[
\boxed{L'=kL},
\]
\[
\boxed{A'=k^2A},
\]
\[
\boxed{V'=k^3V}.
\]

Le coefficient est supérieur à $1$ pour un agrandissement et compris entre $0$ et $1$ pour une réduction.
\end{aretenir}
